\documentclass[10pt,letterpaper]{article}
\usepackage[margin=0.75in]{geometry}
\usepackage{times}
\usepackage[utf8]{inputenc}
\usepackage{amsmath}
\usepackage[colorlinks=true,linkcolor=blue,citecolor=blue,urlcolor=blue]{hyperref}
\usepackage{graphicx}
\usepackage{enumitem}
\usepackage{array}
\usepackage{booktabs}
\usepackage{parskip}
\usepackage{titlesec}

% Reduce spacing
\setlength{\parskip}{6pt}
\setlength{\parindent}{0pt}
\titlespacing*{\section}{0pt}{8pt}{4pt}
\titlespacing*{\subsection}{0pt}{6pt}{3pt}

% Custom section formatting
\titleformat{\section}{\large\bfseries}{\thesection}{1em}{}
\titleformat{\subsection}{\normalsize\bfseries}{\thesubsection}{1em}{}

\begin{document}

\begin{center}
{\LARGE \textbf{Adaptive RDMA Optimization for Distributed LLM\\Training and Inference on Consumer Hardware}}\\[8pt]
{\large Nathan Gopee}\\
M.S. Computer Science, SUNY New Paltz\\
\textit{Proposed Advisors: Dr. Wafi Danesh \& Dr. Ashley Suchy}\\
Target Completion: December 2026
\end{center}

\section{Problem Statement}

Training and serving large language models (7B+ parameters) on consumer GPU clusters is bottlenecked by network communication. Current systems optimize either training (gradient synchronization) or inference (KV cache management) separately, missing opportunities for unified optimization. In our lab with two machines---Cerberus with 2$\times$ RTX 5090s for training, Chimera with 3$\times$ RTX 3090s for inference, connected over 10 Gigabit Ethernet (with potential upgrade to 100GbE)---running a 7B parameter model with 8K context takes 2--3 seconds for Time-To-First-Token (TTFT), a critical latency metric measuring the delay before the first response token arrives \cite{nvidia-nim}. Interactive applications need under 1 second TTFT.

The network has adequate bandwidth. The problem is uniform data treatment: a tiny 4KB attention weight accessed every request waits behind a massive 10MB model layer used once per hundred requests. This head-of-line blocking kills latency despite available bandwidth. Existing solutions require expensive InfiniBand hardware (\$10K+ per node) inaccessible to most academic labs.

\section{Proposed Solution}

I propose a \textbf{unified middleware system} that jointly optimizes both training and inference workloads using adaptive runtime telemetry on commodity Ethernet. By classifying data as ``hot'' (frequently accessed) or ``cold'' (rarely accessed) and routing them appropriately, the system achieves measurable speedups without kernel modifications or enterprise networking equipment.

\textbf{Why This Is Novel:} This would be the first system to jointly optimize gradient communication and KV cache placement using adaptive telemetry. Existing work treats training and inference as separate problems. My approach unifies them under a single classification engine, enabling use cases like continuous self-training where a model simultaneously trains and serves queries.

\textbf{Key Insight:} RDMA (Remote Direct Memory Access) lets machines transfer data directly between memories without CPU involvement, achieving sub-10 microsecond latency versus 50--100 microseconds with TCP/IP. This matters because overhead compounds when moving thousands of tensors per request. By intelligently routing hot data through dedicated RDMA paths while batching and compressing cold data, we eliminate head-of-line blocking.

\section{Technical Approach}

\subsection{Unified Telemetry Monitoring}
The telemetry monitor hooks into vLLM, logging (layer\_id, timestamp, size) for each tensor access in a rolling window with under 1\% overhead. Python decorators wrap vLLM's core operations, maintaining a lock-free circular buffer in shared memory that flushes to Prometheus every 100ms. This provides near-real-time visibility into access patterns for both model weights and KV cache.

\subsection{Adaptive Classification System}
Simple frequency-based classification fails for LLM workloads because access patterns are highly structured: attention weights exhibit phase-dependent bursts during prefill versus decode, KV cache access correlates with conversation context, and gradient tensors follow layer-wise backward propagation order. We use a \textbf{Weighted Finite Automaton (WFA)} that captures these workload-aware patterns.

\textbf{Workload Phase Detection:} The system first identifies the current execution phase $\phi \in \{\text{Prefill}, \text{Decode}, \text{Forward}, \text{Backward}, \text{Idle}\}$ by monitoring operation signatures---prefill shows high parallelism across sequence positions, decode processes single tokens sequentially, backward pass accesses layers in reverse order. Phase detection uses a sliding window of the last 50 operations with majority voting.

\textbf{Weighted Finite Automaton:} For each tensor $t$, we define a WFA $\mathcal{A}_t = (Q, \Sigma, \delta, \lambda, \rho)$ where:
\begin{itemize}[nosep]
\item $Q = \{\text{Cold}, \text{Warm}, \text{Hot}\}$ --- thermal states
\item $\Sigma = \Phi \times \{0,1\}$ --- phase-access pairs (phase $\times$ accessed-or-not)
\item $\delta: Q \times \Sigma \rightarrow Q$ --- transition function
\item $\lambda: Q \times \Sigma \rightarrow \mathbb{R}^+$ --- transition weights learned from access history
\item $\rho: Q \rightarrow [0,1]$ --- routing probability (Hot$\rightarrow$1.0, Warm$\rightarrow$0.5, Cold$\rightarrow$0.0 for dedicated QP)
\end{itemize}

The key insight is that transition weights $\lambda$ are \textit{phase-conditioned}: a tensor might be Hot during Prefill but Cold during Decode (e.g., positional encodings), or Hot during Backward but Cold during Forward (e.g., optimizer states). The WFA learns these conditional weights online using exponential moving averages:
\[
\lambda(q, (\phi, a)) \leftarrow (1-\eta)\lambda(q, (\phi, a)) + \eta \cdot \mathbf{1}[a = 1]
\]
where $\eta = 0.1$ is the learning rate and $\mathbf{1}[a=1]$ indicates an access occurred.

\textbf{Transition Rules:} State transitions fire when accumulated weight crosses adaptive thresholds:
\[
\text{Cold} \xrightarrow{\sum_\phi \lambda(\text{Cold}, (\phi, 1)) > \theta_w} \text{Warm} \xrightarrow{\lambda(\text{Warm}, (\phi_{\text{cur}}, 1)) > \theta_h} \text{Hot}
\]
Thresholds $\theta_w, \theta_h$ adapt based on network congestion (ECN marks increase thresholds, reducing Hot promotions under load). Decay transitions (Hot$\rightarrow$Warm$\rightarrow$Cold) trigger after phase-specific idle timeouts---shorter for Decode phase (500ms) where responsiveness matters, longer for Backward (2s) where batch boundaries create natural gaps.

\textbf{Complexity:} Classification runs in $O(1)$ per tensor: hash lookup for tensor state, constant-time weight update, threshold comparison. Total memory is $O(|T| \cdot |Q| \cdot |\Sigma|)$ for $|T|$ tensors, approximately 360 bytes per tensor.

\subsection{RDMA-Based Smart Routing}
Hot tensors send immediately via dedicated unreliable datagram (UD) queue pairs; cold tensors batch, compress via Top-K sparsification (5\% largest values), and send via shared reliable connection (RC) QPs when batch exceeds 10 tensors or 50ms timeout; warm tensors use shared RC QPs with medium priority. We allocate 4 dedicated QPs for hot traffic per GPU pair, 2 for warm, 1 for cold. Compression uses CUDA kernels for parallel Top-K selection ($<$2\% accuracy impact). GPUDirect RDMA enables zero-copy NIC-to-GPU transfers when supported; otherwise we use pinned host-memory staging.

\subsection{KV Cache Placement \& Migration}
The migration strategy moves KV cache entries to GPUs likely to handle future requests. High-frequency conversations ($>$5 req/min) trigger proactive migration; others fetch on-demand. Protocol: RESIDENT$\rightarrow$MIGRATING$\rightarrow$MIGRATED with atomic pointer swaps after async RDMA copy. We hash the last 128 context tokens to determine target GPU (sticky routing), with hysteresis preventing thrashing: 3 consecutive requests required before migration.

\subsection{Joint Training-Inference Scheduling}
The scheduler dynamically allocates RDMA queue pairs and bandwidth based on workload characteristics, preventing training/inference starvation. Gradient synchronization follows ZeRO patterns with hot/cold classification on frequently-updated layers. Weighted fair queuing gives inference 2$\times$ priority during business hours, training 2$\times$ during off-hours. Periodic bandwidth probes adjust compression ratios to maintain $<$80\% link utilization.

\section{Security Considerations}

RDMA's design prioritizes performance over security, creating vulnerabilities relevant to shared lab environments \cite{bedrock, redmark}. The RDMA credential system (QPN, rKey, PSN) uses predictable values: QPNs are assigned sequentially, rKeys follow linear patterns across memory region registrations, and PSNs have fixed initial values. This enables several attack vectors:

\textbf{Memory Access Attacks:} An attacker who obtains valid credentials can access memory regions belonging to other processes. In shared environments, one user's model weights could theoretically be read by another user's process if credentials leak through side channels or memory disclosure vulnerabilities.

\textbf{Denial of Service:} Malformed packets with incorrect sequence numbers can put queue pairs into error states, disrupting communication. QP exhaustion attacks can deplete NIC resources shared across the system, and bandwidth exhaustion can starve legitimate traffic.

\textbf{Mitigations in This Work:} Since our lab uses dedicated machines for LLM serving (not shared cloud infrastructure), we implement practical mitigations rather than comprehensive security:

\begin{itemize}[nosep]
\item Process isolation via SR-IOV virtual functions (one VF per VM/container)
\item Credential rotation: rKeys regenerated per session, not reused across requests
\item QP limits enforced via cgroup-style resource controls
\item Monitoring: log QP creation/destruction, alert on anomalous patterns
\end{itemize}

For production deployments in multi-tenant environments, network-level defenses using programmable switches (as in Bedrock \cite{bedrock}) or encryption would be necessary. This work focuses on the optimization layer, treating security as a deployment constraint rather than a research contribution.

\section{Implementation Stack \& Related Work}

\textbf{Core Technologies:} pyverbs (RDMA), rdma-core, vLLM 0.6+ (PagedAttention), PyTorch 2.4+, Ray 2.35+, custom C++/pybind11 for zero-copy transfers, CUDA kernels for Top-K compression, Prometheus/Grafana (monitoring).

\textbf{Related Research:} Meta's RDMA deployment \cite{meta-rdma} showed receiver-driven congestion control benefits for ML. vLLM's PagedAttention \cite{vllm} provides efficient KV cache management. Priority-based propagation \cite{p3} and gradient compression \cite{powersgd,accordion} demonstrate value of differentiating data by importance. BytePS \cite{byteps} showed benefits of separating communication from computation.

\textbf{Differentiation:} Existing work optimizes training \cite{meta-rdma,byteps,zero} or inference \cite{vllm,sglang,distserve} separately, often requiring InfiniBand. This middleware jointly optimizes both using runtime telemetry on commodity Ethernet.

\section{Hardware Constraints \& System Requirements}

Everything runs on Ubuntu 24.04 userspace (no kernel mods). Hardware RDMA requires RoCEv2-capable NICs; otherwise we fall back to SoftRoCE with degraded latency.

\textbf{Lab Machines:}
\begin{itemize}[nosep]
\item Cerberus: 2$\times$ RTX 5090 (32GB each), Intel X710 10GbE with hardware iWARP
\item Chimera: 3$\times$ RTX 3090 (24GB each), potential Mellanox ConnectX-6 100GbE upgrade
\end{itemize}

\textbf{Software:} Installation order matters for GPUDirect: CUDA toolkit, then NVIDIA DOCA/MOFED, then nvidia\_peermem module. SR-IOV enables isolated testing via KVM/QEMU with libvirt passthrough. RDMA QPs allocate dynamically per-job, respecting cgroup limits.

\section{Evaluation Methodology}

\textbf{Microbenchmarks:} RDMA latency/bandwidth across message sizes; hot vs cold path comparison; classification overhead; compression throughput.

\textbf{End-to-End:} Inference at varying context lengths and batch sizes; training convergence; combined workloads.

\textbf{Ablations:} Disable classifier, compression, priority routing, and migration individually.

\textbf{Baselines:} vLLM (single-node), vLLM+Ray (multi-node TCP), DeepSpeed ZeRO, NCCL AllReduce.

\section{Possible Timeline}

\begin{table}[h]
\centering
\small
\begin{tabular}{@{}p{1.5cm}p{3cm}p{9.5cm}@{}}
\toprule
\textbf{Months} & \textbf{Phase} & \textbf{Deliverables} \\
\midrule
1--2 & Infrastructure & RoCE/iWARP configuration, vLLM distributed setup, baseline measurements \\
3--4 & RDMA + Compression & Zero-copy library, CUDA Top-K kernels, multi-QP support \\
5--6 & Classification & Access tracking, DFA classifier implementation, monitoring integration \\
7--8 & Integration & Multi-queue scheduler, end-to-end vLLM integration, error handling \\
9--10 & Evaluation & Performance benchmarks, ablation studies, stress testing \\
11--12 & Writing & Thesis draft, revisions, defense preparation, open-source release \\
\bottomrule
\end{tabular}
\end{table}

\section{Technical Contributions}

\textbf{1. Phase-Aware WFA Classifier:} Workload-sensitive tensor classification with learned transition weights and O(1) complexity.

\textbf{2. Multi-Queue RDMA Scheduler:} Priority routing preventing head-of-line blocking with deadlock freedom guarantees.

\textbf{3. Safe Migration Protocol:} Three-state KV cache migration with consistency guarantees.

\textbf{4. End-to-End System:} Practical demonstration on commodity Ethernet showing measurable improvements over baseline approaches.

The broader contribution shows smart software can extract good performance from consumer hardware for LLM workloads, making distributed serving and training accessible to academic labs.

\newpage
\bibliographystyle{plain}
\begin{thebibliography}{99}

\bibitem{nvidia-nim}
NVIDIA Corporation,
\textit{NVIDIA NIM for LLMs: Benchmarking},
2024.
\url{https://docs.nvidia.com/nim/benchmarking/llm/latest/metrics.html}

\bibitem{bedrock}
J. Xing, K. Hsu, Y. Qiu, Z. Yang, H. Liu, A. Chen,
\textit{Bedrock: Programmable Network Support for Secure RDMA Systems},
USENIX Security 2022.

\bibitem{redmark}
B. Rothenberger, K. Taranov, A. Perrig, T. Hoefler,
\textit{ReDMArk: Bypassing RDMA Security Mechanisms},
USENIX Security 2021.

\bibitem{meta-rdma}
A. Gangidi, R. Miao, S. Zheng, et al.,
\textit{RDMA over Ethernet for Distributed AI Training at Meta Scale},
SIGCOMM 2024.

\bibitem{vllm}
W. Kwon, Z. Li, S. Zhuang, et al.,
\textit{Efficient Memory Management for Large Language Model Serving with PagedAttention},
SOSP 2023.

\bibitem{byteps}
Y. Jiang, Y. Zhu, C. Lan, et al.,
\textit{BytePS: A Unified Architecture for Accelerating Distributed DNN Training},
OSDI 2020.

\bibitem{zero}
S. Rajbhandari, J. Rasley, O. Ruwase, Y. He,
\textit{ZeRO: Memory Optimizations Toward Training Trillion Parameter Models},
SC 2020.

\bibitem{p3}
A. Jayarajan, J. Wei, G. Gibson, et al.,
\textit{Priority-based Parameter Propagation for Distributed DNN Training},
MLSys 2019.

\bibitem{powersgd}
T. Vogels, S. Karimireddy, M. Jaggi,
\textit{PowerSGD: Practical Low-Rank Gradient Compression for Distributed Optimization},
NeurIPS 2019.

\bibitem{accordion}
S. Agarwal, H. Wang, K. Lee, et al.,
\textit{Accordion: Adaptive Gradient Communication via Critical Learning Regime Identification},
MLSys 2021.

\bibitem{sglang}
L. Zheng, L. Yin, Z. Xie, et al.,
\textit{SGLang: Efficient Execution of Structured Language Model Programs},
arXiv 2023.

\bibitem{distserve}
Y. Zhong, S. Liu, J. Chen, et al.,
\textit{DistServe: Disaggregating Prefill and Decoding for LLM Serving},
arXiv 2024.

\bibitem{rdma-docs}
Linux RDMA Documentation,
\textit{RDMA Core Userspace Libraries and Daemons}.
\url{https://github.com/linux-rdma/rdma-core}

\end{thebibliography}

\end{document}
